% theme: REFlesson_01nigenkai
\MajorQuestionLesson{連立方程式の意味と解：おかわり}
\SetRowSpace{8mm}

\TypeQuestionDouble{次の数の組を二元一次方程式に代入したとき、等式が成り立つ場合は○、成り立たない場合は×を答えなさい。}{nigenkai}
\QQRow{$x=-3,\ y=5$ は $2x+y=-1$ の解である}{○}{$x=4,\ y=-2$ は $3x-y=13$ の解である}{×}
\QQRow{$x=-5,\ y=-1$ は $x+4y=-8$ の解である}{×}{$x=2,\ y=6$ は $5x-2y=-2$ の解である}{○}
\QQRow{$x=7,\ y=-3$ は $3x+5y=6$ の解である}{○}{$x=-2,\ y=8$ は $4x-y=-15$ の解である}{×}
\QQRow{$x=3,\ y=-5$ は $x-3y=18$ の解である}{○}{$x=-6,\ y=2$ は $3x+4y=-12$ の解である}{×}
\QQRow{$x=-\dfrac{7}{2},\ y=3$ は $2x-5y=-21$ の解である}{×}{$x=\dfrac{1}{2},\ y=-4$ は $4x-3y=14$ の解である}{○}
\QQRow{$x=5,\ y=4$ は $3x-2y=10$ の解である}{×}{$x=-1,\ y=6$ は $5x+2y=7$ の解である}{○}
\QQRow{$x=-\dfrac{3}{2},\ y=-2$ は $4x+5y=-16$ の解である}{○}{$x=8,\ y=1$ は $x-6y=4$ の解である}{×}
\QQRow{$x=6,\ y=-7$ は $4x+3y=5$ の解である}{×}{$x=-8,\ y=5$ は $2x-3y=-31$ の解である}{○}
\QQRow{$x=9,\ y=-4$ は $2x-y=17$ の解である}{×}{$x=-4,\ y=-3$ は $3x-5y=3$ の解である}{○}
\QQRow{$x=-9,\ y=-5$ は $x-4y=11$ の解である}{○}{$x=\dfrac{5}{2},\ y=7$ は $2x+3y=25$ の解である}{×}
